\documentclass[french]{book}
\usepackage{teach}
\usepackage{verbatim}
\usepackage{amsmath,amsfonts,amssymb}
\usepackage{pgfplots}
\usepackage{mwe}
\usepackage{stmaryrd}
\usepackage{yhmath} 
\usepackage{pstricks,pst-plot,pst-text,pst-tree,pst-eucl}
\newcommand{\bsyc}{\left\{\begin{array}{rcl}} % syst�e serr�sur le �al
\newcommand{\esyc}{\end{array}\right.}
\usepackage{picins,picinpar}
\usepackage{pstricks,pst-plot,pst-text,pst-tree,pstricks-add}
\usepackage{pst-eps,pst-fill,pst-node,pst-math}
\usepackage{indentfirst}
%\usepackage[colorlinks=true]{hyperref}
%\usepackage{color}
\newcommand{\ofr}[2]{\dfrac{#1}{#2}}
\newcommand{\arc}[1]{\wideparen{#1}}
\RequirePackage{graphicx}
\RequirePackage{ifpdf}
 \ifpdf
   \DeclareGraphicsRule{*}{mps}{*}{}
 \fi

\newcommand{\sysd}[2]{%
       \left\{
       \begin{aligned}
          #1\\
          #2\\
       \end{aligned}
       \right.%
       }
  \usepackage{fancyhdr}
  \usepackage{tkz-tab}
  \usepackage{amsmath}
  \usepackage{amssymb}
  \usepackage{amsfonts}
  \usepackage{amsthm}
  \usepackage{mathrsfs}

  \usepackage{array}
  \usepackage{multicol}
  \setlength{\multicolsep}{3pt}

  \usepackage{graphicx}
  \graphicspath{{Figures/}}
  \usepackage{pstricks,pst-plot,pst-text,pst-tree,pst-eucl}
  \usepackage[all]{xy}

  \usepackage{fancybox}
  \usepackage{color}

%  \usepackage[frenchb]{babel}
  \usepackage{hyperref}
\usepackage{textcomp}
\usepackage{mathtools}
\usepackage{subfig}
\pgfplotsset{compat=newest}
\newcommand{\I}{i}
\usepackage[all]{xy}
%\usepackage{geometry}
%\geometry{top=1cm, bottom=1cm, left=2cm, right=2cm}
\usepackage[utf8]{inputenc}
\usepackage[french]{babel}
%\usepackage[bottom]{footmisc}
\redefinecolor{thm}{title}{bgcolor}{alizarin}
\redefinecolor{prop}{title}{bgcolor}{meatbrown}
\redefinecolor{defn}{title}{bgcolor}{olivine}
\definecolor{alizarin}{rgb}{0.82, 0.1, 0.26}
\definecolor{meatbrown}{rgb}{0.9, 0.72, 0.23}
\definecolor{olivine}{rgb}{0.6, 0.73, 0.45}
\usepackage{mathrsfs}
%\usepackage{preambule}
%%
\newcommand{\R}{\mathbb {R}}
\newcommand{\bbr}{\mathbb {R}}
\newcommand{\bbc}{\mathbb {C}}
\newcommand{\bbz}{\mathbb {Z}}
%\newcommand{\C}{\mathds {C}}
\newcommand{\N}{\mathbb {N}}
\newcommand{\Z}{\mathbb {Z}}
\newcommand{\Q}{\mathbb {Q}}
\newcommand{\D}{\mathbb {D}}
\newcommand{\HRule}{\rule{\linewidth}{0.5mm}}
%%
\newcommand{\se}{\geq}
\newcommand{\ie}{\leq}
%\newcommand{\imath}{i}
\newcommand{\tm}{\times}
\newcommand{\ab}[1]{\vert #1 \vert}
\newcommand{\nov}[1]{\Vert \overrightarrow{#1} \Vert}
%\newcommand{\ell}{l}
%\newcommand{\ell'}{l'}
%%
\newcommand{\barre}[1]{\overline{#1}}
\newcommand{\ds}{\displaystyle}
\newcommand{\limn}{\lim_{n\rightarrow +\infty}}
\newcommand{\limo}[1][x]{\ds\lim_{#1\rightarrow 0}}
\newcommand{\limite}[2][x]{\ds\lim_{#1\rightarrow #2}}
\newcommand{\limpinf}[1][x]{\ds\lim_{#1\rightarrow +\infty}}
\newcommand{\limminf}[1][x]{\ds\lim_{#1\rightarrow -\infty}}
\newcommand{\limiteg}[2][x]{\ds\lim_{#1\xrightarrow{<}#2}}
\newcommand{\limited}[2][x]{\ds\lim_{#1\xrightarrow{>}#2}}
%%
\newcommand{\vect}[1]{\overrightarrow{#1}}
\newcommand{\oij}{(\textit{O},{\overrightarrow{i}},{\overrightarrow{j}})}
%
\newcommand{\co}[2]{C_#2^#1}
\newcommand{\ud}[1]{\dfrac{#1}{2}}
\newcommand{\cp}[2]{%
        \begin{pmatrix}
        #1\\
        #2
        \end{pmatrix}%
        }
%
\newcommand{\calb}{\mathscr{B}}
\newcommand{\calc}{\mathscr{C}}
\newcommand{\cald}{\mathscr{D}}
\newcommand{\cale}{\mathscr{E}}
\newcommand{\calf}{\mathscr{F}}
\newcommand{\calp}{\mathscr{P}}
\newcommand{\cals}{\mathscr{S}}
\newcommand{\calt}{\mathscr{T}}
\newcommand{\calr}{\mathscr{R}}
\newcommand{\cali}{\mathscr{I}}
\newcommand{\RR}{\calr }
\newcommand{\CR}{\calc}
\newcommand{\IR}{\cali }
\newcommand{\fr}[2]{\dfrac{#1}{#2}}
\newcommand{\ve}[1]{\overrightarrow{#1}}
\newcommand{\pa}[1]{(#1)}
\newcommand{\ps}[2]{\overrightarrow{#1}.\overrightarrow{#2}}
%\newcommand{\bbr}{\mathbb{R}}
\newcommand{\al}{\alpha}
%\newcommand{\la}{\lambda}
\newcommand{\DR}{\cald}
\newcommand{\PR}{\calp}
\newcommand{\cacher}[1]{#1}
\newcommand{\cacheg}[1]{#1}
\newcommand{\cacheb}[1]{#1}
\newcommand{\nor}[1]{ \Vert #1 \Vert}
\newcommand{\abv}[1]{\vert #1 \vert}

\newenvironment{arrayl}[1][rcl]%
	{\setlength{\arraycolsep}{1.5pt}
	 $\begin{array}[t]{#1}}%
	{\end{array}$}%
%\newcommand{\coordpp}[3]{\begin{pmatrix}
%
\newcommand{\suite}[1][u]{\left( #1_{n} \right)_{n\in\N}}
\newcommand{\suitep}[1][u]{\left( #1_{n} \right)_{n\in\N^{*}}}
\usetikzlibrary{arrows}
\newcommand{\ssi}{\Longleftrightarrow}
\newcommand{\poly}[1][a]{#1_0+#1_1x+\cdots+#1_nx^n}
\newcommand{\coordpp}[3]{\begin{pmatrix}
#1\\
#2\\
#3
\end{pmatrix}}
\newcommand{\anglevec}[2]{(\overrightarrow{#1};\overrightarrow{#2})}
\newcommand{\un}{\suite}
\newcommand{\eit}{e^{it}}
\newcommand{\eimt}{e^{imt}}
\newcommand{\vn}{\suite[v]}
%\newcommand{\pa}[1]{(#1)}
\newcommand{\cro}[1]{[#1]}
\newcommand{\abs}[1]{\vert #1 \vert}

\newcommand{\Syst}[2]{\left\{\begin{array}{lcccc} #1 \\ #2 \end{array}\right.}
\renewcommand{\arraystretch}{1.2} 
\newcolumntype{C}[1]{>{\centering\arraybackslash}p{#1cm}}
\newcommand{\Rep}{(O;\ve{i};\ve{j})}
\newcommand{\C}[2]{C_{#1}^{#2}}
\newcommand{\dx}{dx}

\begin{document}
\tableofcontents
\chapter{Matrices}

 \section{Définitions}

\begin{defn}
Une \textbf{matrice} $A$ de dimension $n\times p$ ou de format $(n;p)$ est un tableau de nombres comportant $n$ lignes et $p$ colonnes.
   
   On note $a_{ij}$ l'élément se trouvant à l'intersection de la ligne $i$ et de la colonne $j$.
   
   Lorsque $n=p$ on dit que la matrice est une matrice \textbf{carrée} \textbf{d'ordre} $n$.
   
\end{defn} 

\begin{exemple}
   
    $A=\begin{pmatrix}
    2 & -5 & 4 \\ 
   6 & 3 & -8
   \end{pmatrix} $ est une matrice de format $(2;3)$
   
   Une matrice carrée d'ordre 2 est une matrice de la forme: $\begin{pmatrix}
    a_{11} & a_{12}  \\ 
   a_{21} & a_{22} 
   \end{pmatrix}$
   
\end{exemple}
   
\subsection{Matrices particulières}
   
    Si $n=1$, $A$ est une \textbf{matrice ligne}.
   
     Si $p=1$, $A$ est une \textbf{matrice colonne}.
   
    Si tous les coefficients sont nuls, $A$ est une \textbf{matrice nulle}.
   
\subsection{Égalité de matrices}
   \begin{defn}
   
    Deux matrices sont égales lorsqu'elles ont le même format et les mêmes coefficients aux mêmes emplacements.
   
   \end{defn}
   
\begin{defn}
   
   La matrice \textbf{transposée} d'une matrice $A$ de format $(n;p)$ est la matrice de format $(p;n)$, notée $A^{T}$ obtenue en échangeant les lignes et les colonnes de $A$.
   
\end{defn}   

\begin{exemple}

Si $A=\begin{pmatrix}
    2 & -5 & 4 \\ 
   6 & 3 & -8
   \end{pmatrix} $  alors $A^{T}=\begin{pmatrix}
    2 & 6 \\ 
   -5 & 3 \\
   4 &-8
   \end{pmatrix} $ 
   
\end{exemple}
   
 \subsection{Addition de deux matrices de même format}
   
    On appelle \textbf{somme } de deux matrices de même format la matrice obtenue en additionnant les coefficients de même emplacement.
   
\subsection{Multiplication d'une matrice par un nombre réel}
   
    Pour un réel $k$ et une matrice $A$, on note $kA$ la matrice $M$ dont l'élément $m_{ij}$ est égal à $ka_{ij}$
   
\begin{exemple}
   
   $2\begin{pmatrix}
    2 & -5 & 4 \\ 
   6 & 3 & -8
   \end{pmatrix} = \begin{pmatrix}
    2\times 2 & 2\times (-5) &2\times 4 \\ 
   2\times 6 & 2\times 3 & 2 \times (-8)
   \end{pmatrix}=\begin{pmatrix}
    4 & -10 & 8 \\ 
   12 & 6 & -16
   \end{pmatrix} $ 
   
  \end{exemple}
   
\begin{prop}
   
    $A,B$ et $C$ sont des matrices de même format, $O$ est la matrice nulle de même format, $k$ et $k'$ sont deux nombres réels.


   
   \begin{enumerate}
   \item $A+B=B+A$
   
   \item $A+(B+C)=(A+B)+C$
   
   \item $A+O=O+A=A$
   
   \item $0A=O$ et $1A=A$
   
   \item $(k+k')A=kA+k'A$ et $k(A+B)=kA+kB$ 
   \end{enumerate}
\end{prop}   

\subsection{Multiplication de deux matrices}
   
  \begin{defn} 
   $\blacktriangleright$ Le produit de la matrice ligne $L=\begin{pmatrix}
   a_{1}&a_{2}&...&a_{p}
   \end{pmatrix}$ par la matrice colonne $C= \begin{pmatrix}
   b_{1}\\
   b_{2}\\
   .\\
   .\\
   .\\
   b_{p}
   
   \end{pmatrix}$
      est le nombre $LC=a_{1}b_{1}+a_{2}b_{2}+...+a_{p}b_{p}=\displaystyle\sum_{k=1}^{p}a_{k}b_{k}$
      
   $\blacktriangleright$ Le \textbf{produit} de la matrice $A=\left( a_{ij}\right) $ de format $(n;p)$ par une matrice $B=\left( b_{ij}\right) $ de format $(p;r)$ est la matrice, notée $AB$, de format $(n;r)$ dont le coefficient $(c_{ij})$ est le produit de la matrice ligne $i$ de $A$ par la matrice colonne $j$ de $B$:
   
 \begin{center}
  \boxed{ $$c_{ij}=\displaystyle\sum_{k=1}^{p}a_{ik}b_{kj}$$}
 \end{center}
 \end{defn}
 
\begin{exemple}
 
 $A$ de format $(2;3)$ $B$ de format $(3;2)$, $C=AB$ de format $(2;2)$
  \begin{center}
 \begin{tabular}{cc}
   \begin{pspicture}(2,2)
  \psline{*->}(0.4,-0.6)(2.4,0.6)
  \psline{*->}(1,-0.6)(2.4,0.1)
  \psline{*->}(1.6,-0.6)(2.4,-0.4)
  \end{pspicture}  &$ \begin{pmatrix}
   1 & -1 \\ 
   3 & 2 \\
   -3 & 4 
   \end{pmatrix} $ \\ 
 $\begin{pmatrix}
    1 & 3 & 5 \\ 
   2 & 4 & 6
   \end{pmatrix} $ & $\begin{pmatrix}
    -5 & 25  \\ 
   -4 & 30 
   \end{pmatrix} $ \\ 
 
 \end{tabular} 
 \end{center}
 
 \end{exemple}

\begin{rem}

 
  Si les matrices $AB$ et $BA$ sont définies, $\boxed{ \text{en général } AB\neq BA}$.
  \end{rem}
  
  
\begin{prop}

   $A,B$ et $C$ sont des matrices dont les formats permettent les calculs indiqués, $k$ est un réel
  
  \begin{enumerate}
  \item $A(BC)=(AB)C$
  
  \item $A(B+C)=AB+AC$
  
  \item $(A+B)C=AC+BC$
  
  \item $(kA)B=A(kB)=k(AB)$
\end{enumerate}   

\end{prop}

\subsection{Matrices unités}
\begin{defn}
Soit $n$ un entier naturel non nul. On appelle \textbf{matrice unité d'ordre $n$} la matrice $I$, carrée d'ordre $n$ dont tous les coefficients sont nuls sauf ceux de la diagonale principale (éléments $a_{ii}$) qui sont égaux à 1.
\end{defn}

\begin{exemple}

La matrice unité d'ordre 3 est $I=\begin{pmatrix}
    1 & 0 & 0 \\ 
   0 & 1 & 0 \\
   0 & 0 & 1
   \end{pmatrix} $
  
\end{exemple}

   \section{Inverse d'une matrice carrée}
  
   \begin{defn}
   $A$ est une matrice carrée d'ordre $n$. On dit qu'une matrice $B$, carrée d'ordre $n$, est \textbf{l'inverse} de $A$ si elle vérifie $AB=I$ et $BA=I$.
	\end{defn}   
   
\begin{prop}
   
   Si la matrice carrée $A$ d'ordre $n$ admet une inverse, celle-ci est unique. On la note $A^{-1}$.
   
\end{prop}
   
 \section{Puissance d'une matrice carrée}
   \begin{defn}
   
   $A$ est une matrice carrée et $n \in \mathbb{N}^{*}$. La puissance n-ième de la matrice $A$, notée $A^{n}$, est la matrice définie par: $A^{n}=\underbrace{AA....A}_{\text{n fois}}$
   
    Par convention, $A^{0}=I$
   
  \end{defn}
  
   \begin{thm}
   
   Pour $n \in \mathbb{N}^{*}$ et pour tous nombres réels $a$ et $b$,
    
    \begin{center}
    $\begin{pmatrix}
    a & 0  \\ 
   0 & b 
   \end{pmatrix}^{n}=\begin{pmatrix}
    a^{n} & 0  \\ 
   0 & b^{n} 
   \end{pmatrix}$
    \end{center}
   \end{thm}
   
    \begin{demo}
    par récurrence: exercice
	\end{demo}
	
	    
\begin{thm}
    
    Soit $A=\begin{pmatrix}
    a & b  \\ 
   c & d 
   \end{pmatrix}$ une matrice carrée d'ordre 2.
   
   1) Si $ad-bc\neq 0$, $A$ admet une inverse $\boxed{A^{-1}=\dfrac{1}{ad-bc}\begin{pmatrix}
    d &-b  \\ 
   -c & a 
   \end{pmatrix}}$
   
   2) Si $ad-bc=0$, $A$ n'a pas d'inverse.
   
  \end{thm}
   
\begin{demo}
   
   1) Soit $B=\dfrac{1}{ad-bc}\begin{pmatrix}
    d &-b  \\ 
   -c & a 
   \end{pmatrix}$. On vérifie que $AB=BA=I=\begin{pmatrix}
    1 & 0  \\ 
   0 & 1 
   \end{pmatrix}$
   
   2)  On suppose que $A$ admet une inverse $A'$
   
   Soit $B=\begin{pmatrix}
    -c & a  \\ 
   -c & a 
   \end{pmatrix}$ et $C=\begin{pmatrix}
    d &-b  \\ 
   d & -b 
   \end{pmatrix}$
   
   $BA=\begin{pmatrix}
    -c & a  \\ 
   -c & a 
   \end{pmatrix}$
    $ \begin{pmatrix}
    a & b  \\ 
   c & d 
   \end{pmatrix}=\begin{pmatrix}
    0 & 0  \\ 
   0 & 0
   \end{pmatrix}=O$
   De même, $CA=O$
   
   $B=B(AA')=(BA)A'=O$ donc $a=c=0$
   
   $C=C(AA')=(CA)A'=O$ donc $d=b=0$, $A$ est la matrice $O$ 
   
   ce qui implique que: $AA'=I=O$ ce qui est impossible: $A$ n'a pas d'inverse.\\
  \end{demo} 
   
   \begin{rem}
    Le nombre $ad-bc$ s'appelle le \textbf{déterminant }de la matrice $A$\\
   \end{rem}
   
 \section{Application aux systèmes linéaires}
   
   Tout système linéaire de $n$ équations à $n$ inconnues peut s'écrire sous la forme matricielle $AX=B$ où $A$ est la matrice carrée d'ordre $n$ des coefficients du système, $X$ est la matrice colonne des inconnues et $B$ est la matrice colonne formée par les seconds membres des équations.
   
   Si la matrice carrée $A$ est inversible alors le système a une unique solution égale à $A^{-1}B$.
   
\section{Matrices et suites }
\subsection{Suites de matrices}

Une suite de matrices colonnes de taille $k$ ( $k \in \mathbb{N}, k\geq 2$) est une fonction de $\mathbb{N}$ dans l'ensemble des matrices colonnes de taille $k$.

\begin{defn}

On dit que la suite de matrices colonnes $(X_{n})$ de taille $k$ est \textbf{convergente} si les $k$ suites formées par les termes correspondant à la même ligne sont convergentes.

La limite de la suite est alors la matrice colonne formée des $k$ limites obtenues.
\end{defn}

\subsection{Suites de la forme $U_{n+1}=AU_{n}+B$}

\begin{prop}

Soit une suite de matrices colonnes $(U_{n})$ de taille $k$ telle que, pour tout $n \in \mathbb{N}$,

\begin{center}
$\boxed{U_{n+1}=AU_{n}, \text{ où } A \text{ est une matrice carrée d'ordre }k}$
\end{center}

Alors, pour tout $n$ de $\mathbb{N}$ , $\boxed{U_{n}=A^{n}U_{0}}$\\

\end{prop}

\begin{demo}
\vspace{3cm}
\end{demo}

\begin{prop}


Soit une suite de matrices colonnes $(U_{n})$ de taille $k$ vérifiant pour tout $n \in \mathbb{N}$ , $U_{n+1}=AU_{n}+B$, où $A$ est une matrice carrée non nulle d'ordre $k$ et $B$ une matrice colonne de taille $k$.

S'il existe une matrice $C$ telle que $C=AC+B$ alors le terme général de cette suite peut s'écrire:

\begin{center}
$\boxed{U_{n}=A^{n}(U_{0}-C)+C}$
\end{center}

\end{prop}

\begin{rem}
Si $I-A$ est inversible alors $C=(I-A)^{-1}B$
\end{rem}
 
\begin{demo}
 
 Comme  $U_{n+1}=AU_{n}+B$ et $C=AC+B$, par différence on a: $U_{n+1}-C=A(U_{n}-C)$
 
 La suite $(V_{n})$ telle que $V_{n}=U_{n}-C$ vérifie $V_{n+1}=AV_{n}$ donc $V_{n}=A^{n}V_{0}$
 
 D'où: $U_{n}-C=A^{n}(U_{0}-C)$ soit $U_{n}=A^{n}(U_{0}-C)+C$
 
\end{demo} 

\subsection{Convergence des suites vérifiant $U_{n+1}=AU_{n}+B$}
 
 Soit une suite de matrices colonnes vérifiant $U_{n+1}=AU_{n}+B$
 
 On suppose qu'il existe une matrice $C$ telle que $C=AC+B$
 
 1) Si $U_{0}=C$, la suite converge vers $C$
 
 2) Si $U_{0}\neq C$ et si la suite $(A^{n})$ converge vers une matrice $A'$ alors la suite $(U_{n})$ converge vers
 
 \begin{center}
 $A'(U_{0}-C)+C$
 \end{center}
 
 
\begin{demo}
\vspace{3cm}
\end{demo}
 

\section{Marche Aléatoire}

 \subsection{Marche aléatoire entre deux états}
    \begin{defn}
         
          On considère un système qui n'a que deux états possibles A et B et qui évolue par étapes successives.
           
           On note $p$ la probabilité qu'il passe de A à B.
           
           On note $q$ la probabilité qu'il passe de B à A.
           
           Le \textbf{graphe probabiliste} ci-contre donne l'évolution du système d'une étape à la suivante.
           
           \vspace{0.7cm}
          \begin{center}
           \begin{psmatrix}[mnode=circle]
              A &  & B 
           \end{psmatrix}
           \psset{arrows=->,shortput=nab}
           \nccircle[angleA=0]{1,1}{0.5cm}^{1-p}
           \ncarc[arcangle=15]{1,1}{1,3}^{p}
           \ncarc[arcangle=15]{1,3}{1,1}^{q}
           \nccircle[angleA=0]{1,3}{0.5cm}^{1-q}
          \end{center}
          
          \vspace{0.3cm}
          
          On définit la \textbf{matrice de transition} par: $T=\begin{pmatrix}
    1-p & p  \\ 
   q & 1-q
   \end{pmatrix}$
   
   \end{defn}
   
\begin{rem}
   
   Tous les coefficients appartiennent à $[0;1]$
   
   Pour chaque ligne, la somme des coefficients est 1.
   
   \end{rem}
   
  \begin{defn}
   
   Pour $n \in \mathbb{N}$ , on note:
   
   $A_{n}$ l'événement: "à l'étape $n$ le système est dans l'état A".
   
   $B_{n}$ l'événement: "à l'étape $n$ le système est dans l'état B".
   
   $a_{n}=p(A_{n}), b_{n}=p(B_{n})$. On a: $\boxed{a_{n}+b_{n}=1}$
\end{defn}   
   
\begin{defn}
   
   La matrice \textbf{ligne} $P_{n}=(a_{n}\hspace{0.3cm} b_{n})$ est appelée la \textbf{répartition de probabilité à l'étape $n$}.
   
  \end{defn}

\begin{prop}
   
   Pour $n \in \mathbb{N},\boxed{P_{n+1}=P_{n}T}$
   
   \end{prop}
   
   \begin{demo}
   
   
   \begin{tabular}{cc}
   \psset{nodesep=0mm,levelsep=20mm,treesep=10mm}
\pstree[treemode=R]{\Tdot}
{
\pstree
{\Tdot~[tnpos=a]{$A_{n}$}\taput{\small $a_{n}$}}
{
\Tdot~[tnpos=r]{$A_{n+1}$}\taput{\small $1-p$}
\Tdot~[tnpos=r]{$B_{n+1}$}\tbput{\small $p$}
}
\pstree
{\Tdot~[tnpos=a]{$B_{n}$}\tbput{\small $b_{n}$}}
{
\Tdot~[tnpos=r]{$A_{n+1}$}\taput{\small $q$}
\Tdot~[tnpos=r]{$B_{n+1}$}\tbput{\small $1-q$}
}
}
\begin{tabular}{rl}


 $p(A_{n+1})=$&$ p(A_{n}\cap A_{n+1})+p(B_{n}\cap A_{n+1})$\\
 =&$(1-p)a_{n}+qb_{n}$ \\
 -----------------&-----------------------------------------\\
 $p(B_{n+1})=$&$p(A_{n}\cap B_{n+1})+p(B_{n}\cap B_{n+1})$\\
 =&$pa_{n}+(1-q)b_{n}$
 \end{tabular}\\ 
 
   \end{tabular} 
   
  
\begin{center}
 \begin{tabular}{cc}
    &$ \begin{pmatrix}
    \hspace{1.1cm}1-p & p \textcolor{white}{.} \\ 
  \hspace{1.1cm} q & 1-q \textcolor{white}{.} 
   \end{pmatrix} $ \\ 
 $P_{n}T=\begin{pmatrix}
     a_{n} & b_{n}
   \end{pmatrix} $ & $\begin{pmatrix}
    (1-p)a_{n}+b_{n}q & pa_{n}+(1-q)b_{n}  
   
   \end{pmatrix} $ \\ 
 
 \end{tabular} 
 \end{center}

  On a bien:  $P_{n+1}=P_{n}T$
  
  \end{demo}
  
\begin{prop}

   Pour tout $n \in \mathbb{N}:P_{n}=P_{0}T^{n}$
 \end{prop}
   
  \begin{demo}
  Par récurrence.
\end{demo}   
   
   \begin{defn}
   
   On appelle \textbf{répartition stable de probabilité} une matrice ligne $P$ dont tous les coefficients sont positifs et de somme 1 vérifiant: 
   
   \begin{center}
   $\boxed{P=PT}$
   \end{center}
   
   \end{defn}
   
   \begin{thm}
   Avec les notations précédentes, si $(p;q)\neq (0;0)$ et $(p;q)\neq (1;1)$ alors:
   
   1) il existe une unique répartition stable de probabilité $P$ donnée par:
   
   \begin{center}
   $\boxed{P=\left( \dfrac{q}{p+q}\hspace{0.4cm}\dfrac{p}{p+q}\right) }$
   \end{center}
   
   2) la suite $(P_{n})$ converge vers $P$, indépendamment   de $P_{0}$
   
   \end{thm}
   
\begin{demo}

   $0\leq p \leq 1$ et $0\leq q \leq 1$ donc: $0\leq p+q\leq 2$
   
   Comme $p+q\neq 0$ et $p+q\neq 2$ on a: $0 < p+q<  2$ puis $-2<-p-q<0$ et enfin:
   
   \begin{center}
   $\boxed{-1<1-p-q<1}$ (*)
   \end{center}
   
   On a vu que: $a_{n+1}=(1-p)a_{n}+qb_{n}$ et $b_{n}=1-a_{n}$ donc: $a_{n+1}=(1-p-q)a_{n}+q$
   
   A partir de cette dernière formule, on peut démontrer par récurrence que:
   
   \begin{center}
   $\boxed{a_{n}=(1-p-q)^{n}\left( a_{0}-\dfrac{q}{p+q}\right) +\dfrac{q}{p+q}}$ puis $\boxed{b_{n}=\dfrac{p}{p+q}-(1-p-q)^{n}\left( a_{0}-\dfrac{q}{p+q}\right) }$
   \end{center}
   
   D'après (*), $\lim_{n \to +\infty}\limits a_{n}=\dfrac{q}{p+q}$ et $\lim_{n \to +\infty}\limits b_{n}=\dfrac{p}{p+q}$. On a bien: $\lim_{n \to +\infty}\limits P_{n}=P$
   
   \begin{center}
    \begin{tabular}{cc}
    &$ \begin{pmatrix}
    \hspace{1.1cm}1-p & p \textcolor{white}{.} \\ 
  \hspace{1.1cm} q & 1-q \textcolor{white}{.} 
   \end{pmatrix} $ \\ 
 $PT=P\Leftrightarrow=\begin{pmatrix}
     x & y
   \end{pmatrix} $ & $\begin{pmatrix}
    (1-p)x+qy & px+(1-q)y  
   
   \end{pmatrix}= (x \hspace{0.3cm}y) $ \\
   
 \end{tabular} 
   \end{center}
     
   On obtient le système: $ \left\lbrace
\begin{array}{l}

 x-xp+yq = x\\[0.3cm]
 
  x+y = 1\\[0.3cm] 
\end{array}
\right.  $ qui donne: $ \left\lbrace
\begin{array}{l}

 x=\dfrac{q}{p+q}\\[0.3cm]
 
  y=\dfrac{p}{p+q}\\[0.3cm] 
\end{array}
\right.  $
   
      
   \end{demo}   
      
\section{Marche aléatoire entre plusieurs états}
 
  Les définitions et propriétés précédentes se généralisent à un système qui peut se trouver dans plusieurs états.
  
\begin{thm}
 
 Si la matrice de transition $T$ a une puissance n'ayant aucun coefficient nul, alors:
 
 1) il existe une unique répartition stable de probabilité $P$ telle que $PT=P$
   
      
   2) la suite $(P_{n})$ converge vers $P$, indépendamment   de $P_{0}$ \end{thm}

\end{document}

